\documentclass[12pt]{article}
\usepackage[utf8]{inputenc}
\usepackage[T1]{fontenc}
\usepackage{hyperref}
\title{Middle High German Language Resources for K-12 Curriculum Development}
\author{Compiled by Claude Sonnet 5.5 (claude-sonnet-5-5) for Andr\'as Kornai}
\date{}
\begin{document}
\maketitle

\section{Introduction}
Middle High German (ISO 639-3 gmh) is the historical stage of High German of roughly 1050--1350. It has no native speakers; it is studied through its medieval texts. Searches found scholarly dictionary and corpus resources but no school primers or K-12 programmes.

\section{Primary Educational Resources}
No school textbooks or primers were found in search results.

\section{Linguistic and Reference Materials}
\begin{itemize}
\item Mittelhochdeutsches W\"orterbuch (MWB), Mainz Academy, Trier office. The online and print edition covers 1050--1350, with a lemma list and a digital citation archive, and is linked to the older dictionaries of Benecke--M\"uller--Zarncke and Lexer: \url{https://www.uni-trier.de/en/forschung/mittelhochdeutsches-woerterbuch/start}.
\item Lexer's \emph{Mittelhochdeutsches Taschenw\"orterbuch}: a catalogue record (a scanned edition is reported to be freely available via the Bavarian State Library): \url{https://search.books2ebooks.eu/vufind/Record/bsb_991029117989707356?lng=pl}.
\end{itemize}

\section{Religious and Community Resources}
None found in search results.

\section{Digital and Audio Resources}
\begin{itemize}
\item Reference Corpus of Middle High German (ReM, 1050--1350), Ruhr-Universit\"at Bochum and partners; about 2 million word forms of annotated texts, CC BY-SA 4.0: \url{https://linguistics.rub.de/rem}, access page \url{https://linguistics.rub.de/rem/access/index.en.html}.
\item Uppsala HistCorp distribution (version 1.0, plain text, tokenised and annotated): \url{https://sprakbanken-clarin.lingfil.uu.se/histcorp/readme/german-readme-ReM}.
\end{itemize}

\section{Curriculum Development Context}
The resources found are aimed at researchers and university study. A school curriculum would need simplified readers and glossaries made from the corpus and dictionaries; none were found ready-made. No audio was found in search results.
\end{document}